% §5.4 extended analysis paragraphs
\paragraph{Score overlap and calibration policy.}
Real ShezhenV3 validation pairs exhibit score overlap (minimum clear $Q_{\mathrm{img}}$ below maximum blur $Q_{\mathrm{img}}$) because tongue texture and coating contribute to Laplacian energy even under mild defocus.
We therefore report both cohort medians and Spearman $\rho$ rather than assuming separable clusters.
Routing threshold $\tau$ is set to the midpoint of medians and clamped to $[0.45,0.65]$, matching the phase-2 protocol but replacing synthetic histograms with 572 real clear images and paired blurs.

\paragraph{Implications for deployment.}
Despite overlap, B2 achieves M2$=1.00$ on the test matrix because the closed-loop uplift after resampling ($\Delta Q$ model) combined with real scores still crosses $\tau$ within $K{=}2$ retries for simulated trials.
Operators can inspect \texttt{flags} (\texttt{blur}, \texttt{underexposed}, \texttt{overexposed}) in JSONL logs to diagnose failures without inspecting raw cloud uploads.

\paragraph{Comparison to synthetic phase-2.}
Synthetic calibration yielded $\tau{=}0.505$ with medians $0.819/0.191$; ShezhenV3 yields $\tau{=}0.498$ with medians $0.564/0.431$ and M5 $\rho{=}0.47$ (vs.\ $\rho{\approx}0.75$ synthetic).
Absolute $Q_{\mathrm{img}}$ shifts downward on real texture, but relative routing benefits (B2 vs.\ B0/B1) persist on the test split, supporting external validity of the architecture beyond toy images.
